% Generated draft: requires mathematical and source review.
\documentclass[11pt]{article}
\usepackage[margin=25mm]{geometry}
\usepackage{fontspec}
\setmainfont{Latin Modern Roman}
\usepackage{amsmath,amssymb,mathtools,unicode-math}
\setmathfont{Latin Modern Math}
\usepackage{enumitem,needspace,xurl,hyperref}
\setlength{\parindent}{0pt}
\setlength{\parskip}{5pt}
\setlength{\emergencystretch}{3em}
\newcommand{\N}{\mathbb{N}}
\newcommand{\Z}{\mathbb{Z}}
\newcommand{\Q}{\mathbb{Q}}
\newcommand{\R}{\mathbb{R}}
\newcommand{\C}{\mathbb{C}}
\newcommand{\F}{\mathbb{F}}
\newcommand{\A}{\mathbb{A}}
\newcommand{\PP}{\mathbb{P}}
\newcommand{\E}{\mathbb{E}}
\newcommand{\eps}{\varepsilon}
\newcommand{\dd}{\,\mathrm{d}}
\begin{document}
\section*{UnsolvedMath ID: 2305038}
\textbf{Source problem code:} catalogue entry
\subsection{Definitions and mathematical statement}
Let $\mathbb D=\{z\in\C:|z|<1\}$ and let $S$ conventionally denote the normalized univalent functions on $\mathbb D$: $f(0)=0$, $f'(0)=1$. The source uses $g\prec f$ without defining the symbol. Under the usual analytic-subordination convention it means $g=f\circ\omega$ for a holomorphic self-map $\omega$ of $\mathbb D$ fixing $0$.

\textbf{Requested proof simplification (P. L. Duren).} The record attributes to Tao-Shing Shah the following best-possible assertions under $g\prec f$, $f\in S$, and $g'(0)/f'(0)\in\R$:
\[
 |g(z)|\leq|f(z)|\quad\text{for }|z|\leq\frac{3-\sqrt5}{2},
\]
\[
 |g'(z)|\leq|f'(z)|\quad\text{for }|z|\leq3-\sqrt8.
\]
Here best possible means the respective radius cannot be increased uniformly over the intended class. The task is to find simpler proofs of the reported sharp results, not to establish new larger radii. The source says Shah's proofs are technically involved, and reports that Goluzin gave simpler proofs with worse constants, whose methods \emph{appear} incapable of obtaining the two displayed sharp bounds.

The supplied hypotheses require verification before these can be treated as precise theorems. Under standard subordination, a real derivative ratio permits $\omega(z)=-z$. For instance $f(z)=z/(1-z)^2\in S$ and $g(z)=f(-z)$ have $g'(0)/f'(0)=-1$, while $|g(-r)|>|f(-r)|$ for $0<r<1$. Thus the literal real-ratio hypothesis conflicts with the first displayed claim under that convention. The record may omit a sign/normalization restriction or use a different meaning of $\prec$; no repair is chosen here.

The research summary says the results are established and the item is open only as a simpler-proof request, with no simpler proof located. This accurately distinguishes the source's intended task from an unresolved truth conjecture, but it does not resolve the formulation mismatch.
\subsection{Short English statement}
Find simpler proofs of Shah's two reported sharp domination-radius inequalities for subordinate functions, with radii $(3-\sqrt5)/2$ and $3-\sqrt8$. The printed real derivative-ratio condition conflicts with standard subordination, so the intended hypotheses must first be checked.
\subsection{Sources}
\begin{itemize}
\item Local UnsolvedMath registry, permanent ID 2305038, source code catalogue entry, ``Research Problems in Function Theory --- Problem 5.38''; attribution: P. L. Duren. Supplied status open, classification OPEN-TRIAGE, updated 2026-08-06.
\item Hayman, \emph{Research Problems in Function Theory} (2018), Problem 5.38, \url{https://arxiv.org/abs/1809.07200}. The record reports source-TeX extraction and access on 2026-07-29; the 2018 update was retained for status triage. Current status and rights require review; the record describes publicly available arXiv source TeX.
\item The statement names Tao-Shing Shah for the sharp results and Goluzin for simpler proofs with weaker constants, without complete references. Its labels C5.1 and C5.2 refer to the two bounds but no external bibliographic details are supplied.
\end{itemize}
\end{document}
